\section*{Chapter \ref{ch:g10:layered-ocean} --- The Layered Ocean}

\begin{solution}{exo:g10:layered-ocean:1}
The \omterm{def:g10:layered-ocean:mixed-layer}{mixed layer} is the surface layer of the ocean in which temperature and
salinity hardly change with depth; the \omterm{def:g10:layered-ocean:thermocline}{thermocline} is the layer below it in
which the temperature falls quickly with depth. The wind and the waves stir
the \omterm{def:g10:layered-ocean:mixed-layer}{mixed layer} (helped in winter by the cooling of the surface).
\end{solution}

\begin{solution}{exo:g10:layered-ocean:2}
\omterm{def:g10:layered-ocean:upwelling}{Upwelling} is the rise of cold water from below the surface layer, replacing
surface water pushed away by the wind. The rising water brings mineral
nutrients into the sunlit layer, where the phytoplankton use them to grow;
the phytoplankton feed the rest of the food chain.
\end{solution}

\begin{solution}{exo:g10:layered-ocean:3}
Tropical Pacific: about \qty{29}{\degreeCelsius} at the surface and
\qty{4.5}{\degreeCelsius} at \qty{1000}{m}. Arctic Ocean: about
\qty{-1.4}{\degreeCelsius} at the surface and close to
\qty{0}{\degreeCelsius} at \qty{1000}{m}.
\end{solution}

\begin{solution}{exo:g10:layered-ocean:4}
At \qty{0}{\degreeCelsius} (\qty{1028.1}{kg/m^3} against \qty{1023.3}{kg/m^3}
at \qty{25}{\degreeCelsius}). Seawater: at \qty{4}{\degreeCelsius} pure water
has a density of \qty{1000.0}{kg/m^3}, seawater of salinity 35 about
\qty{1027.8}{kg/m^3} (read on the chart).
\end{solution}

\begin{solution}{exo:g10:layered-ocean:5}
The float drifts at \qty{1000}{m}, sinks to \qty{2000}{m}, then rises to the
surface measuring temperature, salinity (from the conductivity) and pressure
on the way; at the surface it sends its profile by satellite and sinks again.
Floats drift with the currents all over the ocean and work for years, in all
seasons, far from the ship routes.
\end{solution}

\begin{solution}{exo:g10:layered-ocean:6}
$361.84\times10^6 \times 3.682 \approx \qty{1.33e9}{km^3}$, about
\num{1.33} billion cubic kilometres. $100/3682 \approx 0.027$: under
\qty{3}{\percent} of the mean depth.
\end{solution}

\begin{solution}{exo:g10:layered-ocean:7}
$22.7 - 0.5 = \qty{22.2}{\degreeCelsius}$, reached at \qty{15}{m}: the \omterm{def:g10:layered-ocean:mixed-layer}{mixed
layer} is about \qty{15}{m} deep. Between \qty{20}{m} and \qty{50}{m}:
$(21.8 - 17.3)/30 = \qty{0.15}{\degreeCelsius}$ per metre.
\end{solution}

\begin{solution}{exo:g10:layered-ocean:8}
The sun heats only the surface layer, and the seasonal heating and cooling
are spread only through the \omterm{def:g10:layered-ocean:mixed-layer}{mixed layer}: below the \omterm{def:g10:layered-ocean:thermocline}{thermocline} the water is
cut off from them. In winter the surface cools, becomes denser and sinks into
the water below, and the storms stir the sea: the mixing goes deep. In summer
the sun warms a light surface layer that floats on the colder water; the
gentler winds cannot push it down, and the \omterm{def:g10:layered-ocean:mixed-layer}{mixed layer} stays thin.
\end{solution}

\begin{solution}{exo:g10:layered-ocean:9}
$0.78/0.18 \approx 4.3$ degrees near \qty{10}{\degreeCelsius};
$0.78/0.06 = 13$ degrees near \qty{0}{\degreeCelsius}. In the cold polar seas
small changes of salinity matter far more for the density than changes of
temperature: there, salinity controls the layering.
\end{solution}

\begin{solution}{exo:g10:layered-ocean:10}
The temperature is the same: the deeper water is denser because it is
saltier (33.3 against 28.1). The layer between them, where the salinity
changes quickly with depth, is a \omterm{def:g10:layered-ocean:thermocline}{halocline}; here it is also the \omterm{def:g10:layered-ocean:thermocline}{pycnocline}.
\end{solution}

\begin{solution}{exo:g10:layered-ocean:11}
Pushed down, the parcel keeps its temperature and salinity, so its density.
Surrounded by denser water, it is lighter than the water it displaces: the
buoyancy pushes it back up to its own level. A parcel pushed up would be
denser than its new surroundings and sink back. The layers are stable.
\end{solution}

\begin{solution}{exo:g10:layered-ocean:12}
Mass: $0.001 \times 1028.1 = \qty{1.0281}{kg}$. New volume: $1.0281/1023.3
\approx \qty{0.0010047}{m^3}$, about \qty{1004.7}{mL}: it has expanded by
about \qty{4.7}{mL}, \qty{0.47}{\percent}. The layer of \qty{700}{m} would
thicken by $700 \times 0.0047 \approx \qty{3.3}{m}$.
\end{solution}

\begin{solution}{exo:g10:layered-ocean:13}
In the North Atlantic in February. There the top \qty{150}{m} has almost the
same temperature, hence almost the same density: stirring it costs little.
In the tropical Pacific the density rises by $1026.0 - 1021.6 =
\qty{4.4}{kg/m^3}$ between \qty{50}{m} and \qty{200}{m}: a strong
\omterm{def:g10:layered-ocean:thermocline}{pycnocline}, and the storm would have to push light water down into much
denser water.
\end{solution}

\begin{solution}{exo:g10:layered-ocean:14}
High chlorophyll far from coasts: the band along the equator in the Pacific
(and in the Atlantic), where the trade winds drive an \omterm{def:g10:layered-ocean:upwelling}{upwelling}; and the
high latitudes, such as the North Atlantic, where the winter \omterm{def:g10:layered-ocean:mixed-layer}{mixed layer}
reaches deep and brings nutrients up. Very low chlorophyll: the centres of
the subtropical gyres, strongly \omterm{def:g10:layered-ocean:stratification}{stratified}, where nothing brings deep water
up and the sunlit layer has run out of nutrients.
\end{solution}

\begin{solution}{exo:g10:layered-ocean:15}
Sinking: $1000/(6 \times 3600) \approx \qty{0.046}{m/s}$, about
\qty{4.6}{cm/s}. Rising: $2000/(10 \times 3600) \approx \qty{0.056}{m/s}$,
about \qty{5.6}{cm/s}. $36\,000/2 = 18\,000$ measurements, on average
$2000/18\,000 \approx \qty{0.11}{m}$ apart.
\end{solution}

\begin{solution}{pb:g10:layered-ocean:1}
\textbf{1.} The gyre, by $24.0 - 19.0 = \qty{5.0}{\degreeCelsius}$.

\textbf{2.} The gyre (\qty{20.5}{\degree S} against \qty{13.5}{\degree S}).
Farther from the equator, it should be the cooler one: the latitudes alone
cannot explain why it is \qty{5}{\degreeCelsius} warmer.

\textbf{3.} $24.0 - 0.5 = \qty{23.5}{\degreeCelsius}$, reached between
\qty{50}{m} (23.6) and \qty{100}{m} (22.6), at about $50 + 50 \times
0.1/1.0 = \qty{55}{m}$.

\textbf{4.} No: the temperature falls steadily from the surface, by
\qty{1.6}{\degreeCelsius} in the first \qty{25}{m} and \qty{5.3}{\degreeCelsius}
in the first \qty{100}{m}; the \omterm{def:g10:layered-ocean:thermocline}{thermocline} starts at the surface.

\textbf{5.} Between \qty{200}{m} and \qty{300}{m}: $20.3 - 14.4 =
\qty{5.9}{\degreeCelsius}$ per hundred metres.

\textbf{6.} Off Peru, between \qty{50}{m} and \qty{100}{m}: about
\qty{70}{m} by interpolation (the full profile gives \qty{65}{m}). In the
gyre, between \qty{200}{m} and \qty{300}{m}: about $200 + 100 \times 5.3/5.9
\approx \qty{290}{m}$.

\textbf{7.} Yes, at every level of both columns: lighter water lies on denser
water, and both columns are stable.

\textbf{8.} Saltier water should be denser; yet the gyre's surface water is
lighter: \qty{1024.71}{kg/m^3} against \qty{1025.08}{kg/m^3}.

\textbf{9.} The gyre's water is about 1.3 units saltier, which adds about
$1.3 \times 0.8 \approx \qty{1.0}{kg/m^3}$; it is \qty{5}{\degreeCelsius}
warmer, which removes about $5 \times 0.27 \approx \qty{1.4}{kg/m^3}$ at these
temperatures. The temperature wins.

\textbf{10.} $1026.44 - 1025.08 = \qty{1.36}{kg/m^3}$.

\textbf{11.} $1025.30 - 1024.71 = \qty{0.59}{kg/m^3}$.

\textbf{12.} From the surface to \qty{500}{m} the density rises by
$1026.92 - 1024.71 = \qty{2.21}{kg/m^3}$, of which only \qty{0.59}{kg/m^3}
above \qty{200}{m}: most of it, \qty{1.62}{kg/m^3}, lies below
\qty{200}{m}, in the \omterm{def:g10:layered-ocean:thermocline}{thermocline} between \qty{200}{m} and \qty{500}{m}.

\textbf{13.} Off Peru the density rises by $1026.97 - 1025.08 =
\qty{1.89}{kg/m^3}$ from the surface to \qty{500}{m}, of which
\qty{1.36}{kg/m^3} in the top \qty{200}{m}: most of it is near the surface.

\textbf{14.} At \qty{500}{m} the two sites are almost alike
(\qty{7.9}{\degreeCelsius} and \qty{7.1}{\degreeCelsius},
\qty{1026.97}{kg/m^3} and \qty{1026.92}{kg/m^3}): the same cold water lies
beneath both; the difference is how close it comes to the surface.

\textbf{15.} The wind blowing along the shore towards the equator pushes the
surface water out to sea; water from below rises near the coast to take its
place. The cold water of the \omterm{def:g10:layered-ocean:thermocline}{thermocline} is drawn up, so the \omterm{def:g10:layered-ocean:thermocline}{thermocline}
slopes up towards the coast and reaches almost to the surface.

\textbf{16.} Off Peru: water at \qty{13}{\degreeCelsius} to
\qty{15}{\degreeCelsius}, from below the \omterm{def:g10:layered-ocean:thermocline}{thermocline}, lies within the top
\qty{100}{m}, within reach of the wind's stirring and of the light. In the
gyre it lies below \qty{250}{m}, under a warm, light layer that the wind
cannot stir that deep.

\textbf{17.} The \omterm{def:g10:layered-ocean:upwelling}{upwelling} would stop: the surface water would warm in the
sun, a warm light layer would form on top, and the supply of nutrients to
the sunlit layer would dwindle.

\textbf{18.} The coast of Peru is light blue to yellow (much chlorophyll);
the middle of the South Pacific gyre is among the darkest blue regions of the
map (very little).

\textbf{19.} Phytoplankton need nutrients as well as light. In the gyre the
strong \omterm{def:g10:layered-ocean:stratification}{stratification} keeps the deep, nutrient-rich water far below the sunlit
layer, and the nutrients of the surface layer have long been used up.

\textbf{20.} $1.36/0.59 \approx 2.3$: \textbf{the density increase over the
top \qty{200}{m} is about 2.3 times larger off Peru than in the gyre}. The
\omterm{def:g10:layered-ocean:thermocline}{pycnocline}, with the cold, nutrient-rich water beneath it, has been lifted
into the top \qty{200}{m} off Peru, while in the gyre it lies deeper, below
\qty{200}{m}.
\end{solution}
